\section[F. 凸包]{\myboxvi{      \BF{\TT{F}}} 凸包}

\textit{运行时间限制}   2 seconds

\textit{内存使用限制} 1 Gigabyte

\textit{读入文件} \qquad 标准输入

\textit{输出文件} \qquad 标准输出


\begin{quote}
	\textit{我知道既不是凸包也不是凹包的我为什么要哭的。因为我其实是一只豆包。}
\end{quote}

\subsection{题目描述}

给定 $n$ 个一元二次不等式，每个不等式的形式为：

$$
a_i x^2 + b_i x + c_i  >  0
$$

对任意 $1 \le i \le n$ ，保证 $a_i < 0$ 。

请你判断是否存在\BF{整数} $x$，使得这 $n$ 个不等式同时成立。

\subsection{输入格式}

第一行输入一个整数 $T$ ($1 \le T \le 10^4$)，表示测试数据的数量。

每组测试数据第一行输入一个整数 $n$ ($1 \le n \le 2 \times 10^5$)，表示不等式的个数。

随后 $n$ 行，每行输入三个整数 $a_i, b_i, c_i$ ($-10^5 \le a_i \le -1$, $-10^5 \le b_i, c_i \le 10^5$)，表示第 $i$ 个不等式的三个系数。

保证所有测试数据的 $\sum n \le 2 \times 10^5$。

\subsection{输出格式}

对于每组测试数据，若存在整数 $x$ 同时满足所有不等式，输出一行 \TT{YES}，否则输出一行 \TT{NO}。

注意评测时不会区分 \TT{YES} 和 \TT{NO} 的大小写。换言之，当答案是肯定的时候输出 \TT{yes}、\TT{YES}、\TT{Yes}、\TT{YeS} 等都会被认为是正确的。

\subsection{样例 1}

\begin{minipage}{\linewidth}
	\begin{minipage}[t]{0.49\linewidth}
		\begin{lstlisting}[title=\texttt{standard input}]
5
2
-1 0 4
-1 4 -3
3
-1 100000 0
-1 202 10000
-2 400 19998
2
-1 2 0
-1 0 -1
2
-100000 0 100000
-1 0 4
2
-1 2 0
-1 4 -3

\end{lstlisting}
	\end{minipage}
	\begin{minipage}[t]{0.49\linewidth}
		\begin{lstlisting}[title=\texttt{standard output}]
NO
YES
NO
YES
NO
\end{lstlisting}
	\end{minipage}
\end{minipage}


\subsection{样例 1 解释}

对于第 $4$ 组测试数据，当 $x=0$ 时，两个不等式均成立，因此存在整数解。
